problem:  
Let \( (M^6, J, g) \) be a compact **strict nearly Kähler** manifold with SU(3)-structure forms \((\omega, \psi_+, \psi_-)\) satisfying  
\[
d\omega = 3\,\psi_+, \qquad d\psi_- = -2\,\omega^2.
\]  
The associated Riemannian cone \( (C(M), \bar g = dr^2 + r^2g) \) has holonomy contained in \( G_2 \), and serves as the **asymptotic cone** of many known complete \(G_2\)-metrics. Such asymptotically conical (AC) \(G_2\)-manifolds asymptotic to \(C(M)\) have been constructed for certain homogeneous and inhomogeneous nearly Kähler links \(M\), e.g. by Foscolo–Haskins.

**Problem:**  
Fix a compact strict nearly Kähler 6‑manifold \((M,g,J)\) and its corresponding cone \((C(M),\bar g)\).  
Consider the moduli space \(\mathcal{M}_{AC}(C(M))\) of complete torsion‑free \(G_2\) metrics that are asymptotic to this fixed cone \((C(M),\bar g)\) (in an appropriate weighted sense at infinity), modulo diffeomorphisms asymptotic to the identity.

1. Is \(\mathcal{M}_{AC}(C(M))\) **discrete**, or can it contain nontrivial smooth families of non‑isometric \(G_2\) metrics asymptotic to the same NK cone?  
2. Can one describe analytic or topological obstructions (for instance, through the spectrum of the Laplacian on coclosed primitive \((1,1)\)-forms on \(M\), or through a generalized Hitchin functional) that govern the local structure of \(\mathcal{M}_{AC}(C(M))\)?  
3. Might the dimension of \(\mathcal{M}_{AC}(C(M))\) depend on the “type” of the nearly Kähler manifold \(M\) (homogeneous versus inhomogeneous), suggesting a form of rigidity/flexibility dichotomy?

In essence: **Do distinct complete torsion‑free \(G_2\) geometries exist with the same asymptotic nearly Kähler cone, and what mechanisms determine their uniqueness or multiplicity?**

Why is it a "good" problem:  
This problem probes the **uniqueness and deformation theory of asymptotically conical \(G_2\) manifolds** for a fixed nearly Kähler link, a central unsolved question connecting \(G_2\)‑holonomy geometry and 6‑dimensional SU(3)‑structures. It is mathematically precise and logically consistent: it fixes the asymptotic cone (hence the NK structure) and investigates whether multiple complete \(G_2\)-metrics can asymptote to it.  

The problem’s depth lies in uniting **geometric analysis (elliptic PDEs governing torsion‑free \(G_2\) structures)**, **spectral geometry on NK manifolds**, and **moduli theory of special holonomy metrics**. It captures the delicate analytic question of whether infinitesimal deformations of an AC \(G_2\) metric can integrate to genuine global deformations—an issue that remains largely open. Its resolution would illuminate the rigidity or richness of links between nearly Kähler and \(G_2\) geometries, thus serving as a genuine “good problem” that bridges topology, analysis, and differential geometry.